% ===========================================================================
%  Shared preamble for all MDM4U worksheets — input right after \documentclass.
%  (Files beginning with "_" are skipped by mdm4u-worksheet-publish.mjs.)
% ===========================================================================
\usepackage[letterpaper,top=2.2cm,bottom=1.9cm,left=1.6cm,right=1.6cm,headheight=28pt,footskip=24pt]{geometry}
\usepackage{amsmath,amssymb}
\usepackage[T1]{fontenc}
\usepackage[utf8]{inputenc}
\usepackage{lmodern}
\usepackage{xcolor}
\usepackage[most]{tcolorbox}
\usepackage{enumitem}
\usepackage{multicol}
\usepackage{fancyhdr}
\usepackage{array}
\usepackage{tikz}
\usetikzlibrary{calc}
\usepackage{pgfplots}
\pgfplotsset{compat=1.18}

\definecolor{exblue}{HTML}{4A90E2}
\definecolor{exbg}{HTML}{E6F3FF}
\definecolor{qorange}{HTML}{E69138}
\definecolor{qbg}{HTML}{FFF7CC}
\definecolor{keygray}{HTML}{475569}
\definecolor{keybg}{HTML}{F1F5F9}
\definecolor{iagreen}{HTML}{14653B}
\definecolor{titleblue}{HTML}{1D4ED8}
% Rotating palette so each worked example gets its own colour.
\definecolor{ex0}{HTML}{2563A0}\definecolor{ex1}{HTML}{3B7D3B}\definecolor{ex2}{HTML}{8A5A00}
\definecolor{ex3}{HTML}{A3327A}\definecolor{ex4}{HTML}{6D28A3}\definecolor{ex5}{HTML}{B08900}
\definecolor{ex6}{HTML}{0E7490}\definecolor{ex7}{HTML}{9A3412}\definecolor{ex8}{HTML}{15803D}

\pagestyle{fancy}
\fancyhf{}
\fancyhead[L]{\footnotesize NAME: \underline{\hspace{3.4cm}}}
\fancyhead[C]{\textbf{Integration Academy}\\[-2pt]{\scriptsize Math Simplified \textperiodcentered{} Success Amplified}}
\fancyhead[R]{\footnotesize MDM4U}
\fancyfoot[L]{\scriptsize dr.merie@IntegrationAcademy.ca}
\fancyfoot[C]{\scriptsize\thepage}
\fancyfoot[R]{\scriptsize IntegrationAcademy.ca}
\renewcommand{\headrulewidth}{1.2pt}
\renewcommand{\footrulewidth}{0.4pt}
\setlength{\parindent}{0pt}
\setlength{\parskip}{4pt}

% Examples are NOT breakable, so each example + its solution stays on one page.
\newtcolorbox{example}[2]{colframe=#1,colback=#1!7,coltitle=white,colbacktitle=#1,fonttitle=\bfseries,title={#2},arc=3pt,boxrule=1pt}
\newtcolorbox{qbox}[1]{colframe=qorange,colback=qbg,coltitle=white,colbacktitle=qorange,fonttitle=\bfseries,title={#1},arc=3pt,breakable,boxrule=1pt}
\newtcolorbox{keybox}{colframe=keygray,colback=keybg,coltitle=white,colbacktitle=keygray,fonttitle=\bfseries,title={$\checkmark$\ Answer Key},arc=3pt,breakable,boxrule=1pt}
\newtcolorbox{recapbox}{colframe=keygray,colback=keybg,coltitle=white,colbacktitle=keygray,fonttitle=\bfseries,title={Question Recap},arc=3pt,breakable,boxrule=1pt}
\newtcolorbox{ideabox}{colframe=titleblue,colback=titleblue!6,coltitle=white,colbacktitle=titleblue,fonttitle=\bfseries,title={The Big Ideas},arc=3pt,breakable,boxrule=1.2pt}
\newtcolorbox{learnbox}{colframe=iagreen,colback=iagreen!5,coltitle=white,colbacktitle=iagreen,fonttitle=\bfseries,title={Concept in Focus},arc=3pt,breakable,boxrule=1.2pt}
\newcommand{\lh}[1]{\par\smallskip{\bfseries\color{iagreen}#1}\par\nobreak\smallskip}

\newcommand{\soln}{\par\smallskip\textit{\textbf{Solution.}}\ }
\newcommand{\workspace}[1]{\par\vspace{3pt}\fcolorbox{black!30}{white}{\begin{minipage}[t][#1][t]{\dimexpr\linewidth-2\fboxsep\relax}\ \end{minipage}}\par}
\newcommand{\sech}[1]{\par\vspace{8pt}{\Large\bfseries\color{iagreen}#1}\par\nobreak\vspace{2pt}{\color{black!20}\hrule height1pt}\vspace{6pt}}

% A compact coordinate plot sized for an example box: \eplot{xmin}{xmax}{ymin}{ymax}{body}
\newcommand{\eplot}[5]{\begin{center}\begin{tikzpicture}
\begin{axis}[width=6.8cm,height=4.4cm,axis lines=middle,xlabel={\tiny$x$},ylabel={\tiny$y$},
grid=both,grid style={gray!16},xmin=#1,xmax=#2,ymin=#3,ymax=#4,
xtick distance=2,ytick distance=2,tick label style={font=\tiny},axis line style={-},samples=120]
#5
\end{axis}\end{tikzpicture}\end{center}}

% Standard-normal bell with a shaded region between #1 and #2 (in z): \ncurve{from}{to}{extra body}
\newcommand{\ncurve}[3]{\begin{center}\begin{tikzpicture}
\begin{axis}[width=8.6cm,height=4.2cm,axis lines=middle,xlabel={\tiny$z$},ylabel={},
xmin=-3.8,xmax=3.8,ymin=0,ymax=0.45,xtick={-3,-2,-1,0,1,2,3},ytick=\empty,
tick label style={font=\tiny},axis line style={-},samples=160,domain=-3.8:3.8]
\addplot[exblue,very thick]{1/sqrt(2*pi)*exp(-x^2/2)};
\addplot[exblue!22,fill=exblue!22,domain=#1:#2]{1/sqrt(2*pi)*exp(-x^2/2)}\closedcycle;
#3
\end{axis}\end{tikzpicture}\end{center}}

% A bar chart / histogram from (x,height) coordinates: \hist{ymax}{coords}
\newcommand{\hist}[2]{\begin{center}\begin{tikzpicture}
\begin{axis}[width=8.6cm,height=4.6cm,axis lines=left,ybar,bar width=13pt,
ymin=0,ymax=#2, enlarge x limits=0.12,grid=both,grid style={gray!14},
tick label style={font=\tiny},axis line style={-}]
\addplot[draw=exblue,fill=exbg] coordinates {#1};
\end{axis}\end{tikzpicture}\end{center}}

% A box-and-whisker plot on a number line: \boxplot{min}{Q1}{med}{Q3}{max}{axmin}{axmax}
\newcommand{\boxplot}[7]{\begin{center}\begin{tikzpicture}[yscale=0.5]
\draw[gray!40] (#6,0)--(#7,0);
\foreach \t in {#6,...,#7}{\draw[gray!40] (\t,-0.12)--(\t,0.12) node[below=2pt,font=\tiny]{\t};}
\draw[exblue,very thick] (#2,-0.7) rectangle (#4,0.7);
\draw[qorange,very thick] (#3,-0.7)--(#3,0.7);
\draw[exblue,thick] (#1,0)--(#2,0);\draw[exblue,thick] (#4,0)--(#5,0);
\draw[exblue,thick] (#1,-0.4)--(#1,0.4);\draw[exblue,thick] (#5,-0.4)--(#5,0.4);
\end{tikzpicture}\end{center}}

% A scatter plot: \scatter{xmin}{xmax}{ymin}{ymax}{body}
\newcommand{\scatter}[5]{\begin{center}\begin{tikzpicture}
\begin{axis}[width=7.6cm,height=5cm,axis lines=middle,xlabel={\tiny$x$},ylabel={\tiny$y$},
grid=both,grid style={gray!16},xmin=#1,xmax=#2,ymin=#3,ymax=#4,
tick label style={font=\tiny},axis line style={-}]
#5
\end{axis}\end{tikzpicture}\end{center}}
